\section{模块说明}
\texttt{Admin Portal} 是客户端应用背后的运营控制台。本模块的核心不是只浏览后台页面，而是证明后台与移动端使用同一组后端资源：管理员对 \texttt{ETF Catalogue} 和 \texttt{Model Portfolios} 的变更，在刷新移动端后可以得到验证；同时管理员可通过 \texttt{Dashboard}、\texttt{Users} 和 \texttt{User Portfolios} 监控平台状态。
\section{User Story /\allowbreak{} Video 索引}
\begingroup
\small
\begin{longtable}{L{0.18\linewidth}L{0.24\linewidth}L{0.28\linewidth}L{0.16\linewidth}}
\toprule
\textbf{Story} & \textbf{视频主题} & \textbf{核心路径} & \textbf{Video} \\
\midrule
\hyperlink{us-adm-01}{US-ADM-01} & Admin Functional Acceptance & Dashboard + ETF Catalogue + Model Portfolios + operational visibility & \href{https://jjpvro70sief.jp.larksuite.com/wiki/BBmUwLwV0iRRxdkl3S9jV9ZHpLe?from=from\_copylink}{视频} \\
\bottomrule
\end{longtable}
\endgroup
\section{统一术语}
\begingroup
\small
\begin{longtable}{L{0.24\linewidth}L{0.68\linewidth}}
\toprule
\textbf{UI 原文} & \textbf{文档含义} \\
\midrule
\texttt{Admin Portal} & 面向内部管理员的运营控制台 \\
\texttt{Dashboard} & 平台数据与资源概览页 \\
\texttt{ETF Catalogue} & 管理客户端可用 ETF universe 的后台页面 \\
\texttt{Add by Ticker} & 通过 ticker lookup 快速添加 ETF 的流程 \\
\texttt{Add ETF manually} & 由管理员完整填写 ETF 字段的添加流程 \\
\texttt{Model Portfolios} & 管理客户端 Model Portfolio templates 的后台页面 \\
\texttt{Users} & 搜索与查看用户账号、验证、状态和风险画像的页面 \\
\texttt{User Portfolios} & 查看用户所创建 My Portfolio 的运营页面 \\
\texttt{published} /\allowbreak{} \texttt{draft} /\allowbreak{} \texttt{hidden} & Model Portfolio 的发布状态；仅 \texttt{published} 对客户端可见 \\
\texttt{active} /\allowbreak{} \texttt{inactive} & ETF 是否属于客户端 active product universe 的状态 \\
\bottomrule
\end{longtable}
\endgroup
\prdhr
\hypertarget{us-adm-01}{}
\section{User Story 1 - 通过 \texttt{Admin Portal} 维护并核验客户端内容}
\textbf{Story ID:} \texttt{US-ADM-01}
\textbf{用户故事：} 作为管理员，我希望维护 ETF 目录和 Model Portfolio templates，并在移动端验证变更，同时从统一后台查看用户和 portfolio 运营信息。
\textbf{Video:} \href{https://jjpvro70sief.jp.larksuite.com/wiki/BBmUwLwV0iRRxdkl3S9jV9ZHpLe?from=from\_copylink}{US-ADM-01 验收视频}
\subsection{Walkthrough}
\subsubsection{A. 查看 \texttt{Dashboard}}
\begin{enumerate}
\item 登录 \texttt{Admin Portal} 并打开 \texttt{Dashboard}。
\item 展示 total users、ETF catalogue size、user portfolios、structured notes 和 ETF tier distribution 等概览数据，说明这些资源来自与移动端共享的后端。
\end{enumerate}
\screenshotsSingleCompact{assets/screenshots/adm-us1-01-dashboard.png}{ADM-US1-01}{\texttt{Dashboard} 完整概览，包含主要 counters、ETF tier distribution 和最近活动。}
\screenshotDescSingle{\textbf{ADM-US1-01} \textemdash{} \texttt{Dashboard} 完整概览，包含主要 counters、ETF tier distribution 和最近活动。}
\subsubsection{B. 添加 ETF，并在移动端验证}
\begin{enumerate}
\setcounter{enumi}{2}
\item 先在 iOS simulator 搜索一只当前不存在或未启用的测试 ETF，记录找不到该 ETF 的基线状态。
\item 切换到 \texttt{Admin Portal \textgreater{} ETF Catalogue}。展示按 symbol/\allowbreak{}name/\allowbreak{}issuer 搜索，以及按 tier 或 asset class 筛选。
\item 展示 catalogue 中的 exchange、currency、asset class、region、sector、issuer、expense ratio 等字段。
\item 使用 \texttt{Add by Ticker} 或 \texttt{Add ETF manually} 添加测试 ETF。若采用手动流程，填写 symbol、name、exchange、currency、asset class、region、sector、issuer、inception date、expense ratio 和 tier。
\item 返回 simulator，刷新 ETF search/\allowbreak{}list，确认相同 ETF 已出现。
\end{enumerate}
\screenshotsDoubleCompact{assets/screenshots/adm-us1-extra-etf-catalogue.png}{ETF Catalogue 列表与筛选}{\texttt{ETF Catalogue} 列表、search、tier /\allowbreak{} asset class filters，以及 ETF 维护字段。}{assets/screenshots/adm-us1-extra-add-by-ticker-validation.png}{Add by Ticker 重复记录校验}{\texttt{Add by Ticker} 对已存在的 \texttt{VOO} 显示重复记录校验；此图只证明校验行为，不作为成功创建 ETF 的证据。}
\screenshotDescsTwo{\textbf{ETF Catalogue 列表与筛选} \textemdash{} \texttt{ETF Catalogue} 列表、search、tier /\allowbreak{} asset class filters，以及 ETF 维护字段。}{\textbf{Add by Ticker 重复记录校验} \textemdash{} \texttt{Add by Ticker} 对已存在的 \texttt{VOO} 显示重复记录校验；此图只证明校验行为，不作为成功创建 ETF 的证据。}
\subsubsection{C. 停用 ETF，并在移动端验证}
\begin{enumerate}
\setcounter{enumi}{7}
\item 在 \texttt{ETF Catalogue} 选择一只当前可在移动端搜索到的测试 ETF，将其设为 inactive/\allowbreak{}deactivated。
\item 返回 simulator，刷新 ETF search/\allowbreak{}list，确认该 ETF 不再作为 active product 出现。
\end{enumerate}
\subsubsection{D. 修改 \texttt{Model Portfolios}，并在移动端验证}
\begin{enumerate}
\setcounter{enumi}{9}
\item 先在 simulator 打开一个已发布的 Model Portfolio，记录其 title、display order、risk band、summary 和 ETF allocations 基线。
\item 切换到 \texttt{Admin Portal \textgreater{} Model Portfolios}。展示 template key、display order、title、risk band、return range、horizon、summary、reference note、status 和 ETF allocations。
\item 编辑一个容易识别且适合演示的字段，例如 title、display order、risk band、summary 或 allocation。若修改 allocation，确认总权重为 100\%。保存并保持 template 为 \texttt{published}。
\item 返回 simulator，刷新或重新打开 Model Portfolio 页面，确认客户端显示最新配置。
\item 可选补充发布边界：将专用测试 template 改为 \texttt{draft} 或 \texttt{hidden}，刷新客户端确认不再出现；再恢复测试数据，避免影响后续演示。
\end{enumerate}
\screenshotsDoubleCompact{assets/screenshots/adm-us1-extra-model-portfolios-list.png}{Model Portfolios 管理列表}{\texttt{Model Portfolios} 管理列表，展示 template、显示顺序、风险、状态和维护操作入口。}{assets/screenshots/adm-us1-05-model-portfolio-edit.png}{ADM-US1-05}{\texttt{Create Model Portfolio} 表单，显示主要配置字段、ETF allocations 和 \texttt{Total 100.00\%}；发布状态及保存后的客户端结果仍需在视频中连续核验。}
\screenshotDescsTwo{\textbf{Model Portfolios 管理列表} \textemdash{} \texttt{Model Portfolios} 管理列表，展示 template、显示顺序、风险、状态和维护操作入口。}{\textbf{ADM-US1-05} \textemdash{} \texttt{Create Model Portfolio} 表单，显示主要配置字段、ETF allocations 和 \texttt{Total 100.00\%}；发布状态及保存后的客户端结果仍需在视频中连续核验。}
\subsubsection{E. 查看运营页面}
\begin{enumerate}
\setcounter{enumi}{14}
\item 打开 \texttt{Users}，展示 search、role/\allowbreak{}status filters，以及 verification status、account status、risk profile、last login、created date。使用测试用户或对 PII 打码。
\item 打开 \texttt{User Portfolios}，展示 owner、base currency、objective、horizon、holdings count 和 created time。
\end{enumerate}
\screenshotsDoubleCompact{assets/screenshots/adm-us1-extra-users-redacted.png}{Users 运营列表}{\texttt{Users} 运营列表与筛选区域；可识别的 email 已脱敏。}{assets/screenshots/adm-us1-07-user-portfolios.png}{ADM-US1-07}{\texttt{User Portfolios} 运营列表，包含 search /\allowbreak{} filters、owner、base currency、objective、horizon、holdings 和 created time。}
\screenshotDescsTwo{\textbf{Users 运营列表} \textemdash{} \texttt{Users} 运营列表与筛选区域；可识别的 email 已脱敏。}{\textbf{ADM-US1-07} \textemdash{} \texttt{User Portfolios} 运营列表，包含 search /\allowbreak{} filters、owner、base currency、objective、horizon、holdings 和 created time。}
\subsection{Acceptance Criteria}
\begingroup
\small
\begin{longtable}{L{0.18\linewidth}L{0.30\linewidth}L{0.18\linewidth}L{0.12\linewidth}L{0.10\linewidth}}
\toprule
\textbf{ID} & \textbf{Requirement} & \textbf{Reference} & \textbf{Ownership} & \textbf{Priority} \\
\midrule
ADM-US1-AC01 & \texttt{Dashboard} 展示当前后端返回的平台概览，包括 users、ETF catalogue、user portfolios、structured notes 和 ETF tier distribution 中当前可用的指标。 & US-ADM-01 & Integration & Must \\
ADM-US1-AC02 & \texttt{ETF Catalogue} 支持按 symbol、name 或 issuer 搜索，并支持当前实现提供的 tier /\allowbreak{} asset class filters。 & US-ADM-01 & Integration & Must \\
ADM-US1-AC03 & ETF 记录展示 exchange、currency、asset class、region、sector、issuer 和 expense ratio 等维护所需字段。 & US-ADM-01 & Integration & Must \\
ADM-US1-AC04 & 管理员可以通过 \texttt{Add by Ticker} 或 \texttt{Add ETF manually} 创建一条有效 ETF 记录。 & US-ADM-01 & Integration & Must \\
ADM-US1-AC05 & 添加前在移动端不存在的测试 ETF，在后台成功添加后可通过刷新移动端 search/\allowbreak{}list 找到，且 symbol/\allowbreak{}name 与后台一致。 & US-ADM-01 & Integration & Must \\
ADM-US1-AC06 & 管理员可以停用一只 active ETF；停用后刷新移动端，该 ETF 不再出现在 active product universe。 & US-ADM-01 & Integration & Must \\
ADM-US1-AC07 & ETF 添加或停用请求进行中不可重复提交，完成后不会产生重复记录或冲突状态。 & US-ADM-01 & Integration & Must \\
ADM-US1-AC08 & \texttt{Model Portfolios} 支持维护 template key、display order、title、risk band、return range、horizon、summary、reference note、status 和 ETF allocations。 & US-ADM-01 & Integration & Must \\
ADM-US1-AC09 & Model Portfolio allocation 只有在总权重为 100\% 时才能作为有效 template 保存或发布。 & US-ADM-01 & Integration & Must \\
ADM-US1-AC10 & 后台对已发布 Model Portfolio 的 title、order、risk band、summary 或 allocation 所做修改，在刷新或重新打开移动端后得到反映。 & US-ADM-01 & Integration & Must \\
ADM-US1-AC11 & 只有 \texttt{published} Model Portfolio 对客户端可见；\texttt{draft} 或 \texttt{hidden} template 不出现在客户端目录。 & US-ADM-01 & Integration & Must \\
ADM-US1-AC12 & 客户端显示的是当前 Model Portfolio 配置，不使用写死内容或上一次打开的 template 状态。 & US-ADM-01 & Integration & Must \\
ADM-US1-AC13 & \texttt{Users} 支持 search 和 role/\allowbreak{}status filters，并展示当前可用的 verification status、account status、risk profile、last login 和 created date。 & US-ADM-01 & Integration & Must \\
ADM-US1-AC14 & \texttt{User Portfolios} 展示 owner、base currency、objective、horizon、holdings count 和 created time 等当前可用字段。 & US-ADM-01 & Integration & Must \\
ADM-US1-AC15 & \texttt{Users} 和 \texttt{User Portfolios} 用于运营可见性与排查，不在本 story 中修改普通用户持仓。 & US-ADM-01 & Integration & Must \\
ADM-US1-AC16 & 所有会影响客户端的后台变更都通过 simulator 刷新或重新打开页面验证，不能只以后台 Toast 作为成功证据。 & US-ADM-01 & Integration & Must \\
ADM-US1-AC17 & 视频和截图使用测试数据，或对 email、owner 等个人信息进行打码。 & US-ADM-01 & Integration & Must \\
\bottomrule
\end{longtable}
\endgroup
